\documentclass[12pt,a4paper]{article}

%==================== PACKAGES ====================%
\usepackage[a4paper,margin=1in]{geometry}
\usepackage{xcolor}
\usepackage{graphicx}
\usepackage{float}
\usepackage{booktabs}
\usepackage{listings}
\usepackage{tcolorbox}
\usepackage{hyperref}
\usepackage{fancyhdr}
\usepackage{titlesec}
\usepackage{setspace}
\usepackage{longtable}
\usepackage{array}

%==================== COLORS ====================%
\definecolor{mainblue}{RGB}{0,51,102}
\definecolor{lightblue}{RGB}{220,230,242}
\definecolor{codebg}{RGB}{245,245,245}
\definecolor{darkgreen}{RGB}{0,100,0}

%==================== PAGE STYLE ====================%
\pagestyle{fancy}
\fancyhf{}
\fancyhead[L]{Flight Data Analysis Assignment}
\fancyhead[R]{Python Analysis}
\fancyfoot[C]{\thepage}

%==================== HYPERLINK ====================%
\hypersetup{
    colorlinks=true,
    linkcolor=blue,
    urlcolor=darkgreen
}

%==================== TITLE FORMAT ====================%
\titleformat{\section}
{\color{mainblue}\normalfont\Large\bfseries}
{\thesection}{1em}{}

%==================== CODE STYLE ====================%
\lstset{
    backgroundcolor=\color{codebg},
    basicstyle=\ttfamily\small,
    keywordstyle=\color{blue},
    commentstyle=\color{darkgreen},
    stringstyle=\color{red},
    breaklines=true,
    frame=single,
    columns=fullflexible
}

%==================== DOCUMENT ====================%
\begin{document}

%==================== COVER PAGE ====================%
\begin{titlepage}
\begin{center}

\vspace*{2cm}

{\Huge \textbf{\color{mainblue}Flight Data Analysis Using Python}}

\vspace{1cm}

\Large Assignment Report

\vspace{2cm}

\includegraphics[width=0.5\textwidth]{images.jpeg}

\vspace{2cm}

\renewcommand{\arraystretch}{1.15}

\begin{tabular}{|p{4cm}|p{6cm}|}
\hline
\rowcolor{}
\textbf{Course Name} & Programming with R and Python \\ \hline
\textbf{Assignment Title} & Flight Data Analysis Python \\ \hline
\textbf{Prepared By} & Md. Sariful Islam \\ \hline
\textbf{Student ID} & 12310066 \\ \hline
\textbf{Department} & Statistics \\ \hline
\textbf{University} & Begum Rokeya University,Rangpur \\ \hline
\textbf{Submitted To} & Dr. Md. Siddikur Rahman (Associate Professor) \\ \hline
\textbf{Submission Date} & \today \\ \hline
\end{tabular}

\vfill

{\Large \textbf{\color{mainblue}Department of STATISTICS}}

\end{center}
\end{titlepage}

%==================== TABLE OF CONTENTS ====================%
\tableofcontents
\newpage

%==================== INTRODUCTION ====================%
\section{Introduction}

This assignment presents a detailed analysis of airline flight data using Python programming language. The main objective of this project is to analyze coach ticket prices, delays, flight distance, flight duration, first-class prices, and other important variables.

Different statistical methods, visualization techniques, correlation analysis, hypothesis testing, regression analysis, and logistic regression are used in this assignment.

Python libraries used in this project include:

\begin{itemize}
    \item Pandas
    \item Matplotlib
    \item Seaborn
    \item Scikit-learn
    \item SciPy
\end{itemize}

The analysis helps understand relationships among flight variables and provides useful insights regarding airline ticket pricing.

%==================== DATASET DESCRIPTION ====================%
\section{Dataset Description}

The dataset contains airline flight information including:

\begin{itemize}
    \item Coach ticket prices
    \item First-class ticket prices
    \item Flight delays
    \item Flight distance
    \item Flight duration
    \item In-flight services
    \item Weekend and redeye flight information
\end{itemize}

A random sample of 10,000 observations was used for the analysis.

%==================== QUESTION 1 ====================%
\section{Question 1: Coach Ticket Price Analysis}

\subsection{Python Code}

\begin{lstlisting}[language=Python]
import pandas as pd
import matplotlib.pyplot as plt

# Load dataset
flight = pd.read_csv("flight_data.csv")

# Random sample
flight_sample = flight.sample(
    n=10000,
    random_state=66
)

# Summary statistics
print(flight_sample['coach_price'].describe())

# Mean
mean_price = flight_sample[
    'coach_price'
].mean()

print("Average Coach Price:",
      mean_price)

# Histogram
plt.hist(
    flight_sample['coach_price'],
    bins=30
)

plt.xlabel("Coach Ticket Price")
plt.ylabel("Frequency")
plt.title(
"Distribution of Coach Ticket Prices"
)
plt.show()
\end{lstlisting}

\subsection{Output}

\begin{tcolorbox}[colback=lightblue]
count = 10000 \\
mean = 377.57 \\
std = 67.29 \\
min = 94.53 \\
max = 580.85
\end{tcolorbox}

\subsection{Interpretation}

The average coach ticket price is approximately \$378. Most ticket prices are distributed between \$333 and \$427. A ticket price of \$500 is comparatively high but still reasonable for some flights.

%==================== QUESTION 2 ====================%
\section{Question 2: 8-Hour Flight Price Analysis}

\subsection{Python Code}

\begin{lstlisting}[language=Python]
flight_8hr = flight_sample[
    flight_sample['hours'] == 8
]

print(
flight_8hr['coach_price'].describe()
)
\end{lstlisting}

\subsection{Output}

\begin{tcolorbox}[colback=lightblue]
count = 212 \\
mean = 438.76 \\
min = 242.60 \\
max = 561.75
\end{tcolorbox}

\subsection{Interpretation}

Flights lasting 8 hours have a higher average ticket price compared to general flights. Therefore, a \$500 coach ticket appears more reasonable for long-duration flights.

%==================== QUESTION 3 ====================%
\section{Question 3: Flight Delay Distribution}

\subsection{Python Code}

\begin{lstlisting}[language=Python]
plt.hist(
    flight_sample['delay'],
    bins=40
)

plt.xlabel("Delay Time")
plt.ylabel("Frequency")
plt.title(
"Distribution of Flight Delays"
)

plt.show()

print(
flight_sample['delay'].describe()
)
\end{lstlisting}

\subsection{Output}

\begin{tcolorbox}[colback=lightblue]
mean = 13.14 \\
max = 1497
\end{tcolorbox}

\subsection{Interpretation}

Most flights experience small delays. However, the maximum delay is extremely high, indicating the presence of significant outliers.

%==================== QUESTION 4 ====================%
\section{Question 4: Coach Price vs Distance}

\subsection{Python Code}

\begin{lstlisting}[language=Python]
import seaborn as sns

sns.scatterplot(
    x='miles',
    y='coach_price',
    data=flight_sample
)

plt.xlabel("Miles")
plt.ylabel("Coach Price")
plt.title(
"Coach Price vs Flight Distance"
)

plt.show()

print(
flight_sample[
['miles','coach_price']
].corr()
)
\end{lstlisting}

\subsection{Output}

\begin{tcolorbox}[colback=lightblue]
Correlation = 0.372265
\end{tcolorbox}

\subsection{Interpretation}

There is a moderate positive relationship between flight distance and coach ticket prices. Longer flights generally cost more.

%==================== QUESTION 5 ====================%
\section{Question 5: Flight Hours vs Coach Price}

\subsection{Output}

\begin{tcolorbox}[colback=lightblue]
Correlation = 0.366395
\end{tcolorbox}

\subsection{Interpretation}

Longer flight duration slightly increases coach ticket prices.

%==================== QUESTION 6 ====================%
\section{Question 6: Coach Price vs First-Class Price}

\subsection{Output}

\begin{tcolorbox}[colback=lightblue]
Correlation = 0.757325
\end{tcolorbox}

\subsection{Interpretation}

There is a strong positive relationship between coach prices and first-class prices.

%==================== QUESTION 7 ====================%
\section{Question 7: In-Flight Features Impact}

\subsection{Interpretation}

Flights with meals, entertainment, and WiFi generally have higher coach ticket prices.

\begin{itemize}
    \item WiFi increases average ticket price.
    \item Entertainment significantly affects ticket prices.
    \item Meal facilities also increase pricing.
\end{itemize}

%==================== QUESTION 8 ====================%
\section{Question 8: Weekend vs Weekday Flights}

\subsection{Output}

\begin{tcolorbox}[colback=lightblue]
Weekend Coach Price = 406.08 \\
Weekday Coach Price = 323.04
\end{tcolorbox}

\subsection{Interpretation}

Weekend flights are more expensive than weekday flights.

%==================== QUESTION 9 ====================%
\section{Question 9: Redeye Flight Analysis}

\subsection{Output}

\begin{tcolorbox}[colback=lightblue]
Redeye Average = 316.43 \\
Non-Redeye Average = 380.71
\end{tcolorbox}

\subsection{Interpretation}

Redeye flights are generally cheaper than non-redeye flights.

%==================== QUESTION 10 ====================%
\section{Question 10: Advanced Analysis}

%==================== SUMMARY ====================%
\subsection{Summary Statistics}

\begin{tcolorbox}[colback=lightblue]
Average Coach Price = 377.57 \\
Average Delay = 13.14 \\
Average Flight Hours = 3.63
\end{tcolorbox}

%==================== HYPOTHESIS TEST ====================%
\subsection{Hypothesis Testing}

\begin{lstlisting}[language=Python]
from scipy import stats

stats.ttest_ind(
    weekend_yes,
    weekend_no
)
\end{lstlisting}

\subsection{Output}

\begin{tcolorbox}[colback=lightblue]
p-value = 0.0
\end{tcolorbox}

\subsection{Interpretation}

Since the p-value is extremely small, there is a statistically significant difference between weekend and weekday ticket prices.

%==================== REGRESSION ====================%
\subsection{Regression Analysis}

\begin{lstlisting}[language=Python]
from sklearn.linear_model import LinearRegression

X = flight_sample[
 ['miles','hours','delay']
]

y = flight_sample['coach_price']

model = LinearRegression()

model.fit(X,y)

print(model.coef_)
\end{lstlisting}

\subsection{Output}

\begin{tcolorbox}[colback=lightblue]
[0.02785972, -0.90216435, 0.02711225]
\end{tcolorbox}

\subsection{Interpretation}

Flight distance positively impacts coach ticket prices. Delay has a very small effect on pricing.

%==================== LOGISTIC REGRESSION ====================%
\subsection{Logistic Regression}

\begin{lstlisting}[language=Python]
from sklearn.linear_model import LogisticRegression

log_model.fit(X,y)

print(log_model.score(X,y))
\end{lstlisting}

\subsection{Output}

\begin{tcolorbox}[colback=lightblue]
Accuracy = 95.08\%
\end{tcolorbox}

\subsection{Interpretation}

The logistic regression model performs very well in predicting redeye flights.

%==================== CONCLUSION ====================%
\section{Conclusion}

This project successfully analyzed airline flight data using Python. Several statistical and machine learning techniques were applied to understand pricing behavior and flight characteristics.

The analysis revealed that:

\begin{itemize}
    \item Longer flights generally cost more.
    \item Weekend flights are more expensive.
    \item Redeye flights are cheaper.
    \item In-flight services increase ticket prices.
    \item First-class and coach prices are strongly correlated.
\end{itemize}

Overall, Python proved highly effective for data analysis, visualization, and predictive modeling.

%==================== REFERENCES ====================%
\section{References}

\begin{itemize}
    \item Python Documentation
    \item Pandas Documentation
    \item Matplotlib Documentation
    \item Scikit-learn Documentation
\end{itemize}

\end{document}